\documentclass[conference]{IEEEtran}

\usepackage[T1]{fontenc}
\usepackage[utf8]{inputenc}
\usepackage{amsmath,amssymb,amsfonts}
\usepackage{bm}
\usepackage{booktabs}
\usepackage{multirow}
\usepackage{graphicx}
\usepackage{xcolor}
\usepackage{hyperref}
\usepackage{algorithm}
\usepackage{algorithmic}
\usepackage{cite}
\usepackage{array}
\usepackage{url}

\hypersetup{
  colorlinks=true,
  linkcolor=blue!70!black,
  citecolor=blue!70!black,
  urlcolor=blue!70!black
}

\title{LSG: A Lightweight Supervised Gate for Agentic LLM Memory\\
{\large Trading Fixed Threshold Heuristics for a Calibrated,
  Domain-Adaptive Shallow Classifier}}

\author{
  \IEEEauthorblockN{Shlok Chorge}
  \IEEEauthorblockA{
    Department of Information Technology\\
    Vidyalankar Polytechnic and Padmabhushan Vasantdada\\
    Patil College of Engineering (VPPCOE)\\
    University of Mumbai, Mumbai, India\\
    \texttt{shlokchorge2929@gmail.com}
  }
  \and
  \IEEEauthorblockN{Prithibe Majumder}
  \IEEEauthorblockA{\texttt{majumderprithibe@gmail.com}}
}

\begin{document}
\maketitle

\begin{abstract}
Write-time memory gating is a critical yet underexplored
component of deployed LLM agents.
Existing systems rely on a fixed cosine-similarity threshold
to decide what enters long-term memory---a rule that scales
as $O(n)$ with memory size, ignores domain context, and
yields miscalibrated confidence scores.
We introduce \textbf{LSG} (\emph{Lightweight Supervised Gate}),
a shallow tri-model ensemble of Logistic Regression, XGBoost,
and a Multi-Layer Perceptron over frozen 384-dimensional
Sentence-BERT (\texttt{all-MiniLM-L6-v2}) embeddings,
augmented by a per-domain linear adapter
(\textsc{DomainAdapter}).
Gating is reformulated as a supervised binary classification
task ($\textsc{Store}{=}1$ vs.\ $\textsc{Ignore}{=}0$) using
heuristic pseudo-labels derived from Bitext customer-support
and PersonaChat dialogue corpora.
In zero-shot cross-domain transfer (Bitext$\to$PersonaChat),
LSG improves AUPRC from $0.468$ to $0.592$ ($+26.4\%$) and
reduces Expected Calibration Error from $0.233$ to $0.187$
($-19.8\%$) over the similarity baseline, while achieving
constant-time $O(1)$ inference at $0.777$\,ms---an
$11.1{\times}$ speedup at $n{=}1{,}000$ stored items.
The full ensemble trains in under five seconds on commodity
CPU hardware.
Additional diagnostics include in-domain ceiling performance
($F_1{=}0.761$), few-shot adaptation curves
($k\in\{5,20,100\}$), multi-seed variance ($\pm0.008$ AUPRC),
and a false-positive error taxonomy.
\end{abstract}

\begin{IEEEkeywords}
agentic memory, novelty gating, LLM agents, probability
calibration, domain adaptation, lightweight classifiers,
ensemble learning
\end{IEEEkeywords}

\section{Introduction}

Long-horizon conversational agents must retain task-relevant
facts---user preferences, in-flight task state, named
entities, and environmental constraints---within an external
memory store that is retrieved on demand into the active
context window~\cite{packer2023memgpt,mem0_2025,xu2025amem}.
The write-time gate, which decides for each incoming utterance
whether to \textsc{Store} it durably or \textsc{Ignore} it as
ephemeral filler, is therefore a first-order determinant of
both downstream retrieval quality and system efficiency.
Indiscriminate storage inflates the memory index, degrades
retrieval precision, and bloats the LLM context with
non-informative noise such as \emph{``okay''} or
\emph{``sounds good''}.
Excessive conservatism, conversely, causes the system to
discard critical personal facts needed in subsequent turns.

Existing production systems address this trade-off with a
fixed novelty threshold applied to maximum cosine
similarity between the incoming embedding and the stored
memory matrix.
Although this approach is training-free and conceptually
transparent, it carries three fundamental structural
deficiencies:

\begin{enumerate}
  \item \textbf{Content-Type Agnosticism.}
    Cosine similarity captures geometric novelty, not
    informational value.
    A novel social pleasantry (\emph{``I adore rainy days''})
    and a novel durable fact (\emph{``I am allergic to
    peanuts''}) are equidistant from stored vectors, yet
    their long-term utility differs radically.

  \item \textbf{Domain Insensitivity.}
    A single global threshold cannot adapt to shifts in
    the base rate of storeworthy content between domains
    (e.g., entity-dense customer-support dialogues versus
    persona-rich casual conversation).

  \item \textbf{Miscalibration and $O(n)$ Cost.}
    Similarity scores do not correspond to calibrated
    posterior probabilities, and the sequential scan over
    all $n$ memory vectors incurs per-decision cost
    proportional to memory size.
\end{enumerate}

To address all three limitations within a single
computationally frugal framework, this paper introduces
\textbf{LSG (Lightweight Supervised Gate)}, which
reformulates gating as a calibrated binary classification
task over frozen sentence-transformer representations.

\subsection*{Summary of Contributions}

\begin{itemize}
  \item \textbf{Supervised Gating Formulation.}
    We reframe agent memory write gating as a binary
    classification problem over two complementary dialogue
    domains, establishing heuristic pseudo-labels from
    Bitext and PersonaChat corpora.

  \item \textbf{CPU-Native Lightweight Architecture.}
    LSG combines frozen 384-d SBERT features with a
    weighted tri-model ensemble (Logistic Regression,
    XGBoost, MLP) and a per-domain linear adapter,
    training end-to-end in under five seconds on commodity
    CPU hardware.

  \item \textbf{Constant-Time $O(1)$ Inference.}
    LSG achieves a flat per-sample latency of
    $0.777$\,ms, producing an $11.1{\times}$ speedup
    over the $O(n)$ similarity baseline at
    $n{=}1{,}000$.

  \item \textbf{Calibration and Ranking Improvements.}
    Zero-shot cross-domain transfer improves AUPRC from
    $0.468$ to $0.592$ and reduces ECE from $0.233$ to
    $0.187$ relative to the baseline.

  \item \textbf{Comprehensive Empirical Audit.}
    We provide multi-seed variance ($\pm 0.008$~AUPRC),
    95\% bootstrap confidence intervals, bidirectional
    adaptation curves ($k\in\{5,20,100\}$), feature
    importance rankings, embedding-space geometric
    diagnostics, and a false-positive error taxonomy.
\end{itemize}

\section{Related Work}

\subsection{Agentic Memory Systems}
MemGPT~\cite{packer2023memgpt} conceptualizes the LLM context
window as an operating-system paging hierarchy, swapping
content between working memory and long-term storage.
Mem0~\cite{mem0_2025} provides a production-grade key-value
memory layer backed by dense vector retrieval.
A-MEM~\cite{xu2025amem} constructs a dynamically evolving
knowledge graph that applies Zettelkasten-style note linking
and entry superseding.
Crucially, none of these systems treats the initial
write-time \textsc{Store}/\textsc{Ignore} decision as a
calibrated supervised classification problem with explicit
latency constraints---the precise gap filled by LSG.

\subsection{Sentence Representations for Downstream Tasks}
Sentence-BERT~\cite{reimers2019sbert} produces compact,
semantically rich sentence embeddings through contrastive
fine-tuning of transformer encoders.
The paradigm of freezing pre-trained embeddings and training
only a lightweight head is well-established in
data-constrained settings, yielding near-encoder-level
performance at a fraction of the computational cost.

\subsection{Ensemble Classifiers}
XGBoost~\cite{chen2016xgboost} provides rapid, non-linear
gradient-boosted tree ensembles well-suited to
moderate-dimensional input features.
Scikit-learn~\cite{pedregosa2011sklearn} delivers
production-quality implementations of Logistic Regression
and Multi-Layer Perceptrons optimised for CPU execution.
The combination of heterogeneous base learners through
soft-probability averaging has been shown to reduce
variance and improve tail-class calibration relative to
any single constituent model.

\subsection{Calibration}
Expected Calibration Error (ECE)~\cite{guo2017calibration}
quantifies the systematic gap between predicted confidence
and empirical accuracy across equal-width probability bins.
Well-calibrated gating signals are operationally critical
because downstream memory-pipeline components rely on
gate confidence to apply variable-cost storage decisions.

\subsection{Dialogue Corpora}
The Bitext Customer Support
dataset~\cite{bitext2023} supplies slot- and entity-rich
task-oriented utterances spanning diverse support intents.
PersonaChat~\cite{zhang2018personachat} provides
open-domain chit-chat grounded in explicit persona
descriptions, yielding a corpus of first-person factual
statements mixed with generic social exchanges.
The domain mismatch between these two corpora enables
controlled evaluation of cross-domain transfer
and few-shot adaptation.

\section{Problem Formulation}

\subsection{Gating Decision}
Let $\mathcal{U} = \{x_1, x_2, \ldots, x_N\}$ denote a
temporally ordered stream of conversational utterances in
a single agent session.
For each $x_i$, the memory gate $g_\theta$ must produce a
binary output:
\begin{equation}
  y_i =
  \begin{cases}
    1 & \text{(\textsc{Store}: durable fact, entity, preference)} \\
    0 & \text{(\textsc{Ignore}: filler, backchannel, generic query)}
  \end{cases}
\end{equation}

The gate function $g_\theta(x_i) \mapsto \hat{p}_i \in [0,1]$
estimates the posterior $P(y_i{=}1 \mid x_i)$, and a
threshold $\tau \in (0,1)$ yields the hard decision:
\begin{equation}
  \hat{y}_i = \mathbb{I}[\hat{p}_i \geq \tau]
\end{equation}

\subsection{Evaluation Dimensions}
The gate $g_\theta$ is assessed along three orthogonal axes:
\begin{enumerate}
  \item \textbf{Ranking Quality.}
    Threshold-free metrics: ROC-AUC and AUPRC (primary).
  \item \textbf{Decision Quality.}
    Precision $P$, Recall $R$, and $F_1$-score at $\tau{=}0.5$
    and at the calibrated optimal threshold $\tau^*$.
  \item \textbf{Calibration \& Efficiency.}
    ECE (10-bin) and per-sample wall-clock latency $t_{\inf}$
    benchmarked for $n \in \{10, 50, 100, 200, 500, 1000\}$.
\end{enumerate}

\section{Methodology}

\subsection{Heuristic Pseudo-Label Generation}
Ground-truth annotations are replaced by a deterministic
rule-based labeling function $L(x) \mapsto y \in \{0,1\}$
applied uniformly across both corpora:

\begin{itemize}
  \item \textbf{\textsc{Store}=1 (Persona Statements):}
    Regular-expression matching first-person copula and
    predicate patterns, e.g.~\texttt{\^{}i\,(am|have|like|
    love|hate|prefer|work|live|own|\ldots)}.

  \item \textbf{\textsc{Store}=1 (Task Entities):}
    Presence of template slot markers
    (\texttt{\{\{slot\}\}}), named entities, or
    transaction-critical lexemes
    (\emph{order, refund, delivery, payment, reservation}).

  \item \textbf{\textsc{Ignore}=0 (Backchannels):}
    Pattern-match against a curated set of discourse
    particles and generic acknowledgements
    (\emph{ok, yes, no, sure, thanks, hi, bye,
    sounds good, cool}).

  \item \textbf{\textsc{Ignore}=0 (Short Utterances):}
    Any utterance of $\leq 5$ tokens not matched by the
    above positive patterns is conservatively suppressed.
\end{itemize}

\subsection{Feature Representation}
Each utterance $x_i$ is encoded by the frozen
\texttt{all-MiniLM-L6-v2} Sentence-BERT model:
\begin{equation}
  \mathbf{e}_i = \mathrm{SBERT}(x_i) \in \mathbb{R}^{384}
\end{equation}
To prevent OpenMP threading conflicts between
\texttt{sentence-transformers} and XGBoost at inference time,
embedding computation is isolated within a dedicated Python
subprocess.
MD5-keyed disk caching in \texttt{.emb\_cache/} ensures
single-pass computation across repeated experiment runs.

\subsection{Statistical Novelty Baseline}
The baseline mimics production deployment practice: given
memory matrix $\mathbf{M}_{<i} = [\mathbf{m}_1, \ldots,
\mathbf{m}_n]^\top$, it computes a recency-discounted
maximum cosine similarity score:
\begin{equation}
  s_i = \max_{1 \leq j \leq n}
    \left(\frac{\mathbf{e}_i \cdot \mathbf{m}_j}
               {\|\mathbf{e}_i\|\,\|\mathbf{m}_j\|}
    \cdot \gamma^{i-j}\right),\quad \gamma = 0.95
\end{equation}
The gate fires ($\hat{y}{=}1$) if
$s_i < \delta_{\mathrm{novel}} = 0.85$.
This procedure has per-decision time complexity $O(n)$.

\subsection{LSG Ensemble Architecture}
LSG composes three CPU-trainable classifiers, each operating
over the frozen 384-d SBERT features:

\begin{enumerate}
  \item \textbf{Logistic Regression ($\mathbf{m}_{\mathrm{LR}}$):}
    $\ell_2$-regularised linear model
    ($C{=}0.5$, \texttt{max\_iter}$=2000$,
    \texttt{class\_weight=`balanced'}).

  \item \textbf{XGBoost ($\mathbf{m}_{\mathrm{XGB}}$):}
    Gradient-boosted decision tree ensemble
    (\texttt{n\_estimators}$=300$,
    \texttt{max\_depth}$=5$,
    \texttt{learning\_rate}$=0.05$,
    \texttt{subsample}$=0.8$,
    \texttt{colsample\_bytree}$=0.8$).

  \item \textbf{MLP ($\mathbf{m}_{\mathrm{MLP}}$):}
    Fully-connected network $384 \to 256 \to 64 \to 1$
    (\texttt{max\_iter}$=400$, $\alpha{=}10^{-3}$,
    \texttt{early\_stopping=True}).
\end{enumerate}

The ensemble score is produced by fixed-weight soft-voting:
\begin{equation}
  \hat{p}_{\mathrm{raw}} =
    w_{\mathrm{LR}}\,P_{\mathrm{LR}}(1|\mathbf{e}) +
    w_{\mathrm{XGB}}\,P_{\mathrm{XGB}}(1|\mathbf{e}) +
    w_{\mathrm{MLP}}\,P_{\mathrm{MLP}}(1|\mathbf{e})
\end{equation}
with weights $\mathbf{w} = (0.25,\,0.50,\,0.25)$, assigning
the largest share to XGBoost based on held-out calibration
performance.

\subsection{DomainAdapter and Few-Shot Threshold Calibration}
Cross-domain covariate shift is addressed through a
\textsc{DomainAdapter} module that applies a per-domain
affine transformation to input features without retraining
ensemble base models:
\begin{equation}
  \mathbf{e}_{\mathrm{adapt}} = \mathbf{e}\,\mathbf{W}_d^\top + \mathbf{b}_d
\end{equation}
The bias vector $\mathbf{b}_d \in \mathbb{R}^{384}$ is
initialised to the global-minus-domain mean displacement
$(\bar{\mathbf{X}}_{\mathrm{global}} - \bar{\mathbf{X}}_d)$,
while $\mathbf{W}_d \in \mathbb{R}^{384 \times 384}$ is
initialised to $\mathbf{I}_{384}$.
When $k$ target-domain labeled examples are available,
\texttt{DomainAdapter.adapt()} minimises the inverse
centroid-margin objective via gradient descent:
\begin{equation}
  \min_{\mathbf{b}_d,\,\mathbf{W}_d}
  \left\|
    \bar{\mathbf{e}}_{\mathrm{adapt}}^{(1)} -
    \bar{\mathbf{e}}_{\mathrm{adapt}}^{(0)}
  \right\|_2^{-1}
\end{equation}
Following adaptation, a sweep over the precision-recall curve
recalibrates the global decision threshold:
\begin{equation}
  \tau^* = \argmax_{\tau}
  \frac{2 P(\tau)\,R(\tau)}{P(\tau)+R(\tau)+\varepsilon}
\end{equation}

\section{Experimental Setup}

\subsection{Datasets and Domain Partition}
We sample 2\,000 utterances from the Bitext customer-support
corpus~\cite{bitext2023} (task-oriented domain) and 2\,000
utterances from PersonaChat~\cite{zhang2018personachat}
(open-domain), yielding a combined corpus of 4\,000 turns.
Table~\ref{tab:dataset} summarises class distributions.

\begin{table}[h!]
  \centering
  \caption{Dataset Composition and Label Distribution}
  \label{tab:dataset}
  \renewcommand{\arraystretch}{1.15}
  \begin{tabular}{@{}lrrrr@{}}
    \toprule
    \textbf{Domain} & \textbf{Total} & \textbf{Store} & \textbf{Ignore} & \textbf{\% Store} \\
    \midrule
    Task-Oriented & 2\,000 & 1\,000 & 1\,000 & 50.00\% \\
    Open-Domain   & 2\,000 & 1\,073 &   927  & 53.65\% \\
    \midrule
    \textbf{Combined} & \textbf{4\,000} & \textbf{2\,073} & \textbf{1\,927} & \textbf{51.83\%} \\
    \bottomrule
  \end{tabular}
\end{table}

Stratified splits (seed 42) follow the partition scheme below:
\begin{itemize}
  \item $\mathcal{D}_{\mathrm{src\text{-}tr}}$: 1\,600 Bitext samples (80\%).
  \item $\mathcal{D}_{\mathrm{src\text{-}val}}$: 400 Bitext samples (20\%).
  \item $\mathcal{D}_{\mathrm{pc\text{-}pool}}$: 200 PersonaChat few-shot adaptation pool.
  \item $\mathcal{D}_{\mathrm{pc\text{-}te}}$: 1\,800 PersonaChat evaluation test set.
  \item $\mathcal{D}_{\mathrm{pc\text{-}id}}$: 1\,200 train / 800 test PersonaChat in-domain splits.
\end{itemize}

\subsection{Evaluation Metrics}
Models are assessed on five complementary measures:
\begin{enumerate}
  \item Precision ($P$), Recall ($R$), and $F_1$-score on the
    positive class (\textsc{Store}).
  \item ROC-AUC.
  \item AUPRC (primary ranking metric).
  \item ECE with $M{=}10$ equal-width bins:
    \begin{equation}
      \mathrm{ECE} = \sum_{m=1}^{M}
        \frac{|B_m|}{N}\,
        \bigl|\mathrm{acc}(B_m) - \mathrm{conf}(B_m)\bigr|
    \end{equation}
  \item Per-decision wall-clock latency $t_{\inf}$ (ms),
    evaluated at
    $n \in \{10, 50, 100, 200, 500, 1000\}$.
\end{enumerate}

\section{Results and Empirical Analysis}

\subsection{Main Cross-Domain Comparison}
Table~\ref{tab:main} presents the primary evaluation under
zero-shot cross-domain transfer (Bitext-trained, PersonaChat-tested),
alongside the statistical baseline and the in-domain ceiling.

\begin{table*}[t]
  \centering
  \caption{Main Evaluation Results on PersonaChat Test Set}
  \label{tab:main}
  \renewcommand{\arraystretch}{1.2}
  \begin{tabular}{@{}lrrrrrrrr@{}}
    \toprule
    \textbf{Model} & \textbf{Prec.} & \textbf{Recall} & $\bm{F_1}$
      & \textbf{ROC-AUC} & \textbf{AUPRC} & \textbf{ECE}
      & \textbf{Train (s)} & \textbf{Lat. (ms)} \\
    \midrule
    Statistical Baseline
      & 0.5347 & \textbf{1.0000} & 0.6968
      & 0.4149 & 0.4680 & 0.2331 & \textbf{0.00} & 8.614 \\
    LSG Zero-Shot (no adapter)
      & 0.5548 & 0.8160 & 0.6605
      & 0.5830 & 0.5917 & 0.1870 & 4.55 & \textbf{0.777} \\
    LSG Few-Shot ($N{=}100$)
      & 0.5961 & 0.8380 & 0.6964
      & 0.5689 & 0.6040 & 0.2594 & 4.71 & \textbf{0.777} \\
    LSG In-Domain Ceiling
      & \textbf{0.6582} & 0.9023 & \textbf{0.7611}
      & \textbf{0.8142} & \textbf{0.8120} & \textbf{0.0350}
      & 4.31 & \textbf{0.777} \\
    \bottomrule
  \end{tabular}
\end{table*}

\subsubsection{Ranking and Calibration}
LSG zero-shot raises AUPRC from $0.4680$ to $0.5917$
($+26.4\%$ relative gain) and ROC-AUC from $0.4149$ to
$0.5830$, while ECE falls from $0.2331$ to $0.1870$
($-19.8\%$ relative reduction).
These improvements persist across all five random seeds
(multi-seed AUPRC $\sigma = 0.008$) and are confirmed by
non-parametric bootstrap resampling ($B{=}1\,000$):
AUPRC $95\%$~CI $[0.551, 0.613]$;
$F_1$ $95\%$~CI $[0.638, 0.682]$.

\subsubsection{The $F_1$ Paradox}
At the default threshold $\tau{=}0.5$, the statistical
baseline records a nominally higher $F_1$ ($0.6968$
vs.~$0.6605$) by virtue of storing \emph{every} incoming
utterance (Recall $= 1.000$), yielding a precision
($0.5347$) barely above the class prior ($53.65\%$).
This reveals that the baseline functions as a
\emph{pass-through filter} rather than a selective gate.
LSG, by contrast, exercises genuine discrimination,
sacrificing some recall to achieve substantially higher
precision and calibrated confidence.

\subsection{Few-Shot Domain Adaptation}
Table~\ref{tab:fewshot} documents how LSG performance
evolves as the target-domain labeled budget $N$ increases
from zero to the in-domain ceiling.

\begin{table}[h!]
  \centering
  \caption{Few-Shot Adaptation Curve on PersonaChat Test Set}
  \label{tab:fewshot}
  \renewcommand{\arraystretch}{1.2}
  \begin{tabular}{@{}lrrrp{2.4cm}@{}}
    \toprule
    \textbf{Budget $N$} & $\bm{F_1}$ & \textbf{ROC-AUC} & \textbf{ECE} & \textbf{Notes} \\
    \midrule
    0 (Zero-Shot)     & 0.6605 & 0.5830 & 0.1870 & Out-of-box transfer \\
    10                & 0.6689 & 0.5516 & 0.4233 & Adapter init. \\
    25                & 0.6919 & 0.5509 & 0.4471 & Threshold shift \\
    50                & 0.6887 & 0.5584 & \textbf{0.1383} & ECE optimum \\
    100               & 0.6964 & 0.5689 & 0.2594 & Matches baseline $F_1$ \\
    In-Domain Ceiling & \textbf{0.7611} & \textbf{0.8142} & \textbf{0.0350} & Full supervision \\
    \bottomrule
  \end{tabular}
\end{table}

\noindent ECE reaches its minimum at $N{=}50$ ($0.138$),
suggesting that calibration benefits saturate earlier than
discriminative performance.
$F_1$ converges to the baseline value at $N{=}100$ while
maintaining higher selectivity (Precision $0.596$ vs.\
$0.535$).

\subsection{Multi-Seed Variance and Ablation}
Table~\ref{tab:ablation} isolates the contribution of each
ensemble component under five random seeds.

\begin{table*}[t]
  \centering
  \caption{Ablation Study: Mean ($\pm$ Std) Across Five Seeds, Bitext $\to$ PersonaChat Transfer}
  \label{tab:ablation}
  \renewcommand{\arraystretch}{1.2}
  \begin{tabular}{@{}lrrrl@{}}
    \toprule
    \textbf{Variant} & \textbf{$F_1$ (Mean$\pm$Std)} & \textbf{AUPRC (Mean$\pm$Std)} & \textbf{ROC-AUC (Mean$\pm$Std)} & \textbf{Train (s)} \\
    \midrule
    LSG Full Ensemble + Adapter    & $\mathbf{0.660\pm0.005}$ & $\mathbf{0.583\pm0.008}$ & $\mathbf{0.580\pm0.007}$ & 4.31 \\
    LSG Ensemble (No Adapter)      & $0.660\pm0.004$          & $0.583\pm0.007$          & $0.580\pm0.006$          & 4.55 \\
    Logistic Regression Only       & $0.657\pm0.004$          & $0.570\pm0.004$          & $0.556\pm0.002$          & \textbf{0.10} \\
    XGBoost Only                   & $0.642\pm0.006$          & $0.561\pm0.005$          & $0.548\pm0.005$          & 2.15 \\
    MLP Only                       & $0.648\pm0.008$          & $0.565\pm0.009$          & $0.559\pm0.008$          & 1.95 \\
    TF-IDF Vectoriser + Ensemble   & $0.640\pm0.000$          & $0.536\pm0.000$          & $0.538\pm0.000$          & \textbf{0.10} \\
    \bottomrule
  \end{tabular}
\end{table*}

The full ensemble consistently dominates all single-model
ablations, confirming the complementary nature of the three
base learners.
The marginal benefit of the \textsc{DomainAdapter} is modest
in the zero-shot setting, indicating that its primary utility
manifests under explicit few-shot target supervision.

\subsection{Bidirectional Domain Transfer}
Table~\ref{tab:bidir} evaluates adaptation performance
in both directions: Bitext$\to$PersonaChat and
PersonaChat$\to$Bitext.

\begin{table}[h!]
  \centering
  \caption{Bidirectional Few-Shot Adaptation ($F_1$ and AUPRC)}
  \label{tab:bidir}
  \renewcommand{\arraystretch}{1.2}
  \begin{tabular}{@{}lrrrrl@{}}
    \toprule
    \textbf{Direction} & $\bm{k}$ & $\bm{F_1}$ & $\bm{\Delta F_1}$ & \textbf{AUPRC} & $\bm{\Delta\text{AUPRC}}$ \\
    \midrule
    \multirow{4}{*}{\shortstack[l]{Bitext $\to$\\PersonaChat}}
      & 0   & 0.667          & {---}      & 0.579          & {---}      \\
      & 5   & 0.611          & $-0.056$   & 0.577          & $-0.002$   \\
      & 20  & 0.693          & $+0.026$   & 0.600          & $+0.021$   \\
      & 100 & \textbf{0.695} & $+0.028$   & \textbf{0.604} & $+0.025$   \\
    \midrule
    \multirow{4}{*}{\shortstack[l]{PersonaChat $\to$\\Bitext}}
      & 0   & 0.585          & {---}      & 0.632          & {---}      \\
      & 5   & 0.585          & $+0.000$   & 0.632          & $+0.000$   \\
      & 20  & 0.611          & $+0.026$   & 0.655          & $+0.023$   \\
      & 100 & \textbf{0.611} & $+0.026$   & \textbf{0.656} & $+0.024$   \\
    \bottomrule
  \end{tabular}
\end{table}

The initial regression at $k{=}5$ in the Bitext$\to$PersonaChat
direction is a common pattern in adapter-based few-shot
learning: the adapter overspecialises on very limited
examples before recovering at $k{=}20$.
Both directions show monotone AUPRC improvements from
$k{=}20$ onward, validating the adapter design across
asymmetric domain difficulties.

\subsection{Inference Latency and Memory Scaling}
Table~\ref{tab:latency} records wall-clock per-decision
latency as the stored memory index $n$ grows from 10 to
1\,000 items.

\begin{table}[h!]
  \centering
  \caption{Per-Decision Inference Latency vs.\ Memory Size}
  \label{tab:latency}
  \renewcommand{\arraystretch}{1.2}
  \begin{tabular}{@{}rrrr@{}}
    \toprule
    $\bm{n}$ & \textbf{LSG (ms)} & \textbf{Baseline (ms)} & \textbf{Speedup} \\
    \midrule
    10      & \textbf{0.777} & 1.544 & $2.0\times$ \\
    50      & \textbf{0.777} & 1.966 & $2.5\times$ \\
    100     & \textbf{0.777} & 1.828 & $2.4\times$ \\
    200     & \textbf{0.777} & 3.402 & $4.4\times$ \\
    500     & \textbf{0.777} & 3.658 & $4.7\times$ \\
    1\,000  & \textbf{0.777} & 8.614 & $\mathbf{11.1\times}$ \\
    \bottomrule
  \end{tabular}
\end{table}

LSG latency is invariant to $n$, confirming its $O(1)$
complexity class.
The baseline's latency grows super-linearly beyond $n{=}200$,
consistent with memory-bandwidth effects in large cosine
scans.

\subsection{Feature Importance and Embedding Geometry}
Table~\ref{tab:geometry} presents intrinsic geometric
diagnostics of the frozen SBERT embedding space on the
PersonaChat test set.

\begin{table}[h!]
  \centering
  \caption{Embedding-Space Geometry and Linear Separability}
  \label{tab:geometry}
  \renewcommand{\arraystretch}{1.2}
  \begin{tabular}{@{}lrp{3.2cm}@{}}
    \toprule
    \textbf{Metric} & \textbf{Value} & \textbf{Interpretation} \\
    \midrule
    Within-\textsc{Store} cosine dist.  & 0.1654  & Moderate intra-class cohesion \\
    Within-\textsc{Ignore} cosine dist. & 0.1465  & Slightly tighter clustering \\
    Across-class cosine dist.           & 0.1382  & Significant inter-class overlap \\
    Cosine separation ratio             & 1.1286  & Weak but positive separation \\
    $L_2$ centroid distance             & 0.1430  & Small Euclidean centroid gap \\
    LDA Fisher ratio                    & 0.0115  & Non-zero linear signal \\
    Linear probe accuracy               & 80.78\% & Upper bound on separability \\
    \bottomrule
  \end{tabular}
\end{table}

The cosine separation ratio of $1.129$ confirms that the
SBERT embedding space provides a weak but exploitable
linear signal.
The large gap between the raw across-class distance
($0.138$) and the linear probe accuracy ($80.78\%$) suggests
that the decision boundary is achievable only through
supervised optimisation, not unsupervised distance thresholding.

Logistic Regression feature-coefficient analysis reveals that
the 10 highest-magnitude dimensions are: 238 ($|w|{=}2.029$),
301 ($1.711$), 330 ($1.545$), 84 ($1.495$), 62 ($1.490$),
137 ($1.436$), 80 ($1.364$), 14 ($1.353$), 144 ($1.344$),
and 17 ($1.340$).
The broad distribution of discriminative weight across
non-adjacent dimensions corroborates that storeworthiness is
encoded as a distributed, non-axis-aligned semantic property.

\subsection{Calibration Audit}
Table~\ref{tab:calib} decomposes LSG and baseline calibration
across equal-width confidence bins on the zero-shot test set
($N{=}1\,800$).

\begin{table*}[t]
  \centering
  \caption{Bin-Level Calibration Audit --- LSG Zero-Shot ($N{=}1\,800$)}
  \label{tab:calib}
  \renewcommand{\arraystretch}{1.2}
  \begin{tabular}{@{}crrrrp{3.8cm}@{}}
    \toprule
    \textbf{Bin} & $\bm{n}$ & \textbf{Mean Conf.} & \textbf{Actual Acc.} & \textbf{Gap} & \textbf{Status} \\
    \midrule
    $[0.0,\;0.1)$ &   3 & 0.0866 & 0.3333 & $-0.247$ & Under-confident \\
    $[0.1,\;0.2)$ &  46 & 0.1562 & 0.4130 & $-0.257$ & Under-confident \\
    $[0.2,\;0.3)$ &  83 & 0.2573 & 0.4699 & $-0.213$ & Under-confident \\
    $[0.3,\;0.4)$ & 109 & 0.3495 & 0.5505 & $-0.201$ & Under-confident \\
    $[0.4,\;0.5)$ & 144 & 0.4512 & 0.4028 & $+0.049$ & Over-confident (non-monotone) \\
    $[0.5,\;0.6)$ & 183 & 0.5536 & 0.3934 & $+0.160$ & Over-confident \\
    $[0.6,\;0.7)$ & 262 & 0.6524 & 0.4809 & $+0.172$ & Over-confident \\
    $[0.7,\;0.8)$ & 406 & 0.7525 & 0.5961 & $+0.156$ & Over-confident \\
    $[0.8,\;0.9)$ & 465 & 0.8467 & 0.6065 & $+0.240$ & Over-confident \\
    $[0.9,\;1.0]$ &  99 & 0.9194 & 0.6364 & $+0.283$ & Over-confident \\
    \bottomrule
  \end{tabular}
\end{table*}

LSG exhibits a systematic calibration pattern: the
low-confidence region ($\hat{p} < 0.4$) is under-confident,
while the mid-to-high range ($\hat{p} \geq 0.4$) is
over-confident.
The non-monotone transition at $[0.4, 0.5)$ suggests a
decision-boundary artefact from weighted ensemble probability
averaging.
Post-hoc temperature scaling or Platt scaling targeting the
$[0.5, 0.9]$ range is identified as the most impactful
intervention for future deployment hardening.

\subsubsection{False-Positive Taxonomy}
Of the 640 false positives recorded on the zero-shot test
set, qualitative categorisation identifies three dominant
failure modes:
\begin{enumerate}
  \item \textbf{Longer Ambiguous} (493 cases,
    $\bar{c}{=}0.726$): Extended utterances containing
    incidental location or entity references without
    durable informational intent
    (e.g., \emph{``I was thinking about visiting Paris
    someday''}).
  \item \textbf{Short Question} (108 cases,
    $\bar{c}{=}0.748$): Conversational queries containing
    surface-level nouns classified as entity-bearing
    (e.g., \emph{``What is your favourite book?''}).
  \item \textbf{Short Other} (29 cases,
    $\bar{c}{=}0.763$): Brief social assertions with
    inherently ambiguous store-worthiness.
\end{enumerate}

\section{Discussion and Limitations}

\subsection{Core Empirical Claims}
The experimental results substantiate three principal claims:
\begin{enumerate}
  \item \textbf{Selective Gating Superiority.}
    LSG eliminates the baseline's degenerate pass-through
    behaviour (Recall\,$=1.000$, $F_1$\,boosted trivially),
    delivering a genuine selective gate with $+26.4\%$
    relative AUPRC improvement.
  \item \textbf{Computational Scalability.}
    Replacing $O(n)$ cosine scans with a single $O(1)$
    forward pass yields an $11.1{\times}$ wall-clock
    speedup at $n{=}1{,}000$ stored items.
  \item \textbf{Training Efficiency.}
    The full tri-model ensemble converges in $<5$\,seconds
    on standard laptop CPU hardware, requiring neither
    GPU acceleration nor reinforcement learning reward
    modelling.
\end{enumerate}

\subsection{Limitations and Future Work}

\paragraph{Heuristic Label Noise.}
Pseudo-label quality is bounded by regular-expression
coverage.
Social pleasantries containing incidental entity terms
(e.g., location names in wishes rather than factual claims)
consistently evade the ignoring rules, injecting label noise
that constrains the supervised accuracy ceiling.
Crowdsourced annotation or LLM-assisted labelling would
reduce this noise floor.

\paragraph{Mid-Range Overconfidence.}
The calibration audit confirms systematic overconfidence
for $\hat{p} \geq 0.5$.
Planned mitigations include temperature scaling, Platt
sigmoid recalibration, and isotonic regression, all of
which can be applied post-hoc without ensemble retraining.

\paragraph{Single-Modality Features.}
LSG relies exclusively on sentence-level embeddings, ignoring
turn-level dialogue history, speaker-role signals, and
document-structure cues.
Augmenting features with conversational context embeddings
or lightweight dialogue-act labels is a promising direction.

\section{Conclusion}

This paper presented \textbf{LSG (Lightweight Supervised Gate)},
a shallow ensemble classifier for agentic LLM memory gating
that addresses three fundamental limitations of prevailing
fixed-threshold heuristics: content-type agnosticism, domain
insensitivity, and calibration deficiency.
By combining Logistic Regression, XGBoost, and an MLP over
frozen Sentence-BERT embeddings---augmented with a
parameter-efficient per-domain linear adapter---LSG delivers
statistically significant improvements in cross-domain
ranking quality (AUPRC $0.592$ vs.\ $0.468$, $+26.4\%$),
reduced calibration error (ECE $0.187$ vs.\ $0.233$,
$-19.8\%$), and constant-time $O(1)$ inference at
$0.777$\,ms per decision.
The system trains in under five seconds on commodity CPU
hardware and requires no GPU resources, making it a
practically deployable module for production conversational
AI pipelines.
Future iterations will target label-quality improvement,
post-hoc calibration, and contextual feature enrichment
to close the remaining gap toward the in-domain ceiling
($F_1{=}0.761$, AUPRC $0.812$).

\bibliographystyle{IEEEtran}
\begin{thebibliography}{9}

\bibitem{packer2023memgpt}
C. Packer, S. Wooders, K. Lin, V. Fang, S. G. Patil,
I. Stoica, and J. E. Gonzalez,
``MemGPT: Towards LLMs as Operating Systems,''
\emph{arXiv preprint arXiv:2310.08560}, 2023.

\bibitem{mem0_2025}
Mem0 Team,
``Mem0: Building Production-Ready AI Agents with
Scalable Long-Term Memory,''
2025.

\bibitem{xu2025amem}
W. Xu \emph{et al.},
``A-MEM: Agentic Memory for LLM Agents,''
\emph{arXiv preprint arXiv:2502.12110}, 2025.

\bibitem{reimers2019sbert}
N. Reimers and I. Gurevych,
``Sentence-BERT: Sentence Embeddings using Siamese
BERT-Networks,''
in \emph{Proc. EMNLP-IJCNLP}, 2019.

\bibitem{chen2016xgboost}
T. Chen and C. Guestrin,
``XGBoost: A Scalable Tree Boosting System,''
in \emph{Proc. ACM SIGKDD}, 2016.

\bibitem{pedregosa2011sklearn}
F. Pedregosa \emph{et al.},
``Scikit-learn: Machine Learning in Python,''
\emph{J. Mach. Learn. Res.}, vol.~12, pp.~2825--2830,
2011.

\bibitem{guo2017calibration}
C. Guo, G. Pleiss, Y. Sun, and K. Q. Weinberger,
``On Calibration of Modern Neural Networks,''
in \emph{Proc. ICML}, 2017.

\bibitem{bitext2023}
Bitext,
``Bitext Customer Support LLM Chatbot Training Dataset,''
HuggingFace Datasets,
\texttt{bitext/Bitext-customer-support-llm-chatbot-training-dataset}, 2023.

\bibitem{zhang2018personachat}
S. Zhang, E. Dinan, J. Urbanek, A. Szlam, D. Kiela,
and J. Weston,
``Personalizing Dialogue Agents: I have a dog, do you
have pets too?,''
in \emph{Proc. ACL}, 2018.

\end{thebibliography}

\end{document}
